\renewcommand{\thetable}{3.A.3b}
\begin{tablalafrox}{Requisitos propios ofertados. Fuente: elaboración propia a partir de la Tabla 3.A.3a, los requerimientos de Distribuidora Puelche y el Anexo 2.2 del Subdocumento 2 (LafroX, 2026a).}%
  {tab:3-A-3b}%
  {F{0.062} F{0.099} Y F{0.179} F{0.224} F{0.098} F{0.093}}%
  {\cab{ID} & \cab{Actor} & \cab{Descripción} & \cab{Precondición} & \cab{Resultado esperado} & \cab{Prioridad / Etapa} & \cab{Origen}}
  RF-02.11 & Jefa de Bodega / Gerente de Operaciones & El sistema debe registrar en cada plataforma de cross-docking la lectura de las unidades logísticas al llegar el camión de línea y al cargarlas en el camión de reparto, con la hora de cada evento, y las excepciones (unidad faltante, sobrante o dañada) con su causa. & La transferencia está registrada como stock en tránsito desde Talca (RF-02.06) y las unidades llevan su identificador logístico. & Por cada plataforma y día se obtiene el tiempo real entre consolidación, llegada, desconsolidación y despacho, y el detalle de excepciones, para medir la brecha de 16 horas (LafroX, 2026a, Anexo 2.2, SP-02). & Alta / E1 & Caso, sección 2.3 y cap. 3; SD2, Anexo 2.1, F-05. \\
  RF-06.14 & Jefe de TI / Gerente de Finanzas / Tripulación propia y peoneta / Conductores externos & El sistema debe enviar al ERP, para cada guía de despacho electrónica, la evidencia de recepción capturada en la entrega (firma, fotografía, nombre del receptor o confirmación QR) y recibir de vuelta el estado del acuse de recibo registrado por el ERP, conciliándolo con la guía. & La entrega tiene evidencia capturada (RF-06.11 o RF-06.12) y la guía fue emitida por el ERP. Sin señal, la evidencia queda en cola y se envía al sincronizar. & Cada guía queda con su estado de acuse (recibido conforme, con observaciones o sin acuse) visible el mismo día; las guías sin acuse generan una alerta para gestión. La solución no emite documentos tributarios. & Alta / E1 & SD2, Anexo 2.1, F-09. \\
  RF-06.06 & Tripulación propia y peoneta / Conductores externos & El sistema debe registrar la hora real de llegada y de término de cada entrega. & La entrega está en curso en la aplicación de reparto (RF-06.01). & Cada entrega tiene hora real de llegada y término, que alimenta el cálculo de OTIF (RF-11.03). & Alta / E1 & . \\
  RF-09.10 & Jefa de Calidad & El sistema debe permitir parametrizar por tipo de producto la regla de excursión térmica (desviación y duración). & Los tipos de producto y sus rangos están definidos por Calidad (V-11). & La alerta y la retención de RF-09.05 y RF-09.06 aplican el parámetro vigente, y cada cambio queda registrado. & Media / E1 & Caso, sección 16.1, decisión 4; SD2, Anexo 2.2, S-04. \\
  RF-10.01 & Jefe de abastecimiento & El sistema debe sugerir la reposición de compras por SKU y proveedor a partir de la historia de ventas, con excepciones revisables. & Existe historia de ventas en el ERP (V-09, S-25). & La sugerencia reemplaza la planilla de reposición y queda disponible para revisión (D-05). & Media / E1 & Caso, sección 4.1 y cap. 5. \\
  RF-10.02 & Jefe de abastecimiento & El sistema debe incorporar al cálculo de reposición las promociones comprometidas. & Existe una sugerencia de reposición (RF-10.01) y promociones comprometidas registradas. & La sugerencia incluye las promociones sin agregarlas a mano. & Media / E1 & Caso, sección 4.1. \\
  RF-10.03 & Jefe de abastecimiento & El sistema debe traspasar la sugerencia aprobada al módulo de compras del ERP. & La sugerencia fue aprobada (RF-10.01). & El ERP recibe la sugerencia aprobada y conserva la emisión de la orden de compra. & Media / E1 & Caso, sección 4.1; cap. 10. \\
\end{tablalafrox}
